\section{A measurement case: tags, strings, and uses}\label{sec:evidence}

A corpus query can identify a useful class without recovering the intended
linguistic bearers. The Corpus of Global Web-Based English (GloWbE)
contains country-labelled web sections \citep{daviesfuchs2015}.
An interjection tag, the string \mention{ha}, and a conventional
interjectional use of \mention{ha} select different sets. Confusing them
can turn a retrieval convention into apparent evidence about a category.

The \mention{ha} record makes the consequence concrete. Some tagged
occurrences abbreviate \term{hectare}, especially in agricultural prose.
A regex screen of exported contexts changes Kenya's rank from fifth
to fifteenth or sixteenth under six treatments of laughter and unclear
cases; pooling gives ranks from fourteenth to sixteenth. The large
movement is stable across those choices. Its exact size remains
conditional on the screen and on the exported samples.

The screen was tuned against forty Kenyan and Tanzanian examples;
its resulting 39/40 binary agreement is a development check. It doesn't
establish equivalent accuracy in every country section. The observed
geographic distribution therefore depends on both the linguistic
phenomenon and the measurement procedure. An analysis of category
variation would need independently identified uses and comparable
opportunities for their occurrence.

The retained frequency model answers a narrower question. It compares
prediction of an omitted country--item aggregate while retaining the
same item in nineteen other countries and seventeen other items in
that country. Adding country intercepts has an imprecisely estimated
advantage of 17.9 expected log predictive density units (standard error
13.3). This is prediction within the observed design; new items,
new countries, and the proposed linguistic groupings weren't tested.

Similarly, changing frequency ranks in the Corpus of Historical American
English (COHA) describe a tagged inventory \citep{davies2010}.
They don't establish persistence of a semantic or grammatical profile.
Two further analyses, of derivatives and \mention{fie}, fail their
intended confirmatory interpretation: the primary derivative pipeline
assigned semantic residue rather than measuring it, and the \mention{fie}
input retained ineligible optical character recognition substitutions.
Their confirmatory claims are withdrawn.

The \href{inquiry-relative-evidence-supplement.pdf}{evidence supplement}
preserves the complete records: GloWbE measurement and model comparisons
(Section~A), the derivative audit (Section~B), the \mention{fie} audit
(Section~C), and inventory turnover (Section~D). These records locate
the failure at the relevant claim. A useful retrieval class survives
as a tool, but its output cannot establish a linguistic category
without evidence connecting that class to independently identified uses.

\section{The general consequence}\label{sec:discussion}

Named groupings expose what a broad interjection judgment can obscure.
The grammatical analysis of prenominal \mention{damn} concerns its
integration into the noun phrase; the expressive analysis concerns
its contribution to interpretation. Depictive \mention{Bang!} poses
another semantic question. Continuers and repair initiators organize
different sequential environments. These relations needn't fit a
fixed matrix with one column for each discipline.

\subsection{Many-to-many relations and local failure}

Dupr\'e's pluralism and Khalidi's crosscutting kinds make room for
multiple defensible groupings. The linguistic cases add a practical
requirement: specify which relation a proposed grouping preserves and
which objects participate in it. Haspelmath's division between descriptive
categories and comparative concepts supplies a related lesson about
analytical purpose. Here, the same care is needed within a language
and within an inquiry.

The expressive--depictive comparison establishes that several groupings
can be useful for one question about meaning contribution. A procedural
account asks a different question about encoding, potentially cutting
across those groupings. If two encoding accounts conflict, however,
the conflict remains. Neither putting them in different fields nor
assigning each a new label makes their claims compatible.

Failure has correspondingly local consequences. A grouping can lack
reliable application conditions, obscure the stated descriptive contrast,
or fail to support a projected property. A historical explanation can
lack independently supported source evidence. Revising the grouping or
the question creates a new proposal; it doesn't preserve the support
of a failed claim. Other distinctions survive only where they have
their own evidence.

This discipline also limits positive conclusions. The semantic comparison
favours preserving expressive and depictive contributions, while leaving
open whether separate classes, subtypes, or dimensions best serve a
further task. The phonological comparison rejects an unsupported transfer
to grammatical membership. The \mention{gee} history supports a particular
explanation. These are substantive outcomes without a claim that this
paper has validated a complete family of empirical categories.

\subsection{Overlap, mechanisms, and category standing}

For common bearers, intersections are well-defined and may be useful.
Their value depends on the task and alternatives, including representations
of separate memberships and their dependencies. No automatically privileged
category follows from an overlap among grammatical interjections,
expressives, and interactional resources. Nor is a useful intersection
excluded merely because its membership can be derived.

Distinct categories needn't be statistically independent, and they can
share causal processes. Strong correlation doesn't make two questions
identical; a shared history doesn't establish every present property.
A causal network is one possible account, provided its variables and
links have independent support. Hill's considerations help assess
specified causal interpretations; they don't turn covariation into
a mechanism by being applied to a diagram.

Gradience is another separate issue. A class can have graded membership,
sharp subtypes, or unresolved cases due to missing evidence. A prototype
account organizes typicality and similarity \citep{rosch1975,taylor2003}.
The assessment here asks what such organization accomplishes and what
it licenses an investigator to infer. Either a prototype or a categorical
representation can serve a task; neither gains empirical category standing
from its form alone.

A field therefore supplies no fixed quota of categories, and an inquiry
needn't end with just one. A grouping may serve several questions across
fields, yet warrant only specified projections under specified conditions.
The historical label \term{interjection} identifies a meeting point for
these inquiries. The evidence, rather than the label or departmental
boundary, determines which proposed categories deserve recognition.

\section{Conclusion}\label{sec:conclusion}

English interjections expose a general difference between assembling a
class and establishing an empirical category. Descriptive classes can
make contrasts visible; category claims require co-occurring properties
that support induction for particular purposes. Explanations require
further evidence about how an outcome arose or is maintained. These
achievements should be assessed for named groupings, not allocated in
advance to grammar, semantics, and pragmatics.

The expressive--depictive comparison shows why one semantic inquiry can
need several groupings. The phonological case shows why evidence for
one grouping leaves a different membership question open. A history
can explain an outcome without validating an ancestry category, and a
corpus tag can aid retrieval without identifying the intended sign-types.
In each case, making the bearer, grouping basis, and intended contribution
explicit changes what would count as support or counterevidence.
An undifferentiated category judgment conceals those differences; the
account developed here makes them available for adjudication.

\ifblind\else
\section*{Acknowledgements}

I compiled empirical material for the Wikipedia article \enquote{English
interjections}, of which I'm the primary author. This paper reanalyses
that material as a worked case for inquiry-relative category assessment.
The archived data and analysis code are available at
\url{https://github.com/BrettRey/interjections-hpc/tree/master}.
The linked evidence supplement distinguishes descriptive results from
confirmatory interpretations withdrawn after audit.
\fi

\clearpage
\begingroup
\setlength{\emergencystretch}{3em}
\printbibliography
\endgroup
\end{document}
